\subsection{When is the twin a valid proxy?}
For passive, image-based reconstruction, the error budget is dominated by the reconstruction itself
(view coverage, texture, appearance), and the twin reproduces it: on average the real and simulated CD of
NeRF, 3DGS and MVS differ by $-0.02$ to $+0.32$~mm, with the same ranking in both domains. A twin with
calibrated cameras and replayed poses can therefore be used to select and tune these methods. For active
depth, the twin models the sensor as an ideal range measurement plus random noise, so it captures the
precision of the L515 but not its trueness: the depth-to-color extrinsics, a range bias that depends on the
material, and subsurface scattering are absent. The result is an optimistic prediction by a factor of ten.
Closing this gap requires a measured error model of the sensor (per-camera rigid bias and
material-dependent range bias), which the paired setup makes directly measurable: the per-method ICP
translation is an estimate of that bias.

\subsection{Uncertainty of the comparison}
Several error sources limit the resolution with which gaps can be resolved. The CAD-to-print deviation is
\printmae~mm MAE (max \printmax~mm). The ChArUco board poses have a median detector error of about 0.4~mm
in the twin, where the true pose is known, and the twin reproduces the board projections within 0.11~px.
The 4-DOF registration of the CAD to the real reconstructions has an RMS residual of about 1~mm on the
inlier points. Differences below roughly 0.5~mm between the real and simulated CD should therefore not be
interpreted as significant, which is the case for MVS and 3DGS on average.

\subsection{Choice of metric}
Accuracy, completeness and F-score lead to different conclusions than CD alone. On the real captures,
MVS has the best mean median accuracy of the image-based methods (0.86~mm versus 0.93~mm for NeRF) but a
worse completeness (1.34 versus 1.00~mm) and the largest H95 among them (10.8~mm), because it leaves
low-texture regions empty; its CD is therefore similar to that of 3DGS, whose clouds are complete but
thicker (median accuracy 1.22~mm). Reporting both directions of the distance and several thresholds is
necessary, in agreement with \cite{zielinski2026comparative,knapitsch2017tanks}.

\subsection{Limitations}
This study uses one capture per object, so the repeatability of the real measurements, and thus confidence
intervals of the gap, are not available. Materials and lighting of the twin were matched visually rather
than measured, the objects were not placed at their real pose on the board (the twin uses a canonical
placement, Section~\ref{sec:registro}), so the lighting direction relative to the object differs between
domains, the depth noise model is provisional, and the CAD-to-print deviation was measured with a
caliper on one artifact only. Color fidelity, which is relevant for the radiance-field methods, was not
evaluated. Finally, five objects and one trajectory do not cover the variability of textures,
reflectances and numbers of views that a full factorial study would require.
